% !TEX root = ../diplomarbeit.tex

\chapter{Projektmanagement}

% --------------------------------------------------------
% Projektplanung und Meilensteine
% --------------------------------------------------------
\section{Planung}

\emph{Dieser Abschnitt dokumentiert die Planung der Diplomarbeit, einschließlich der definierten Meilensteine, Zeitlinie und Aufgabenzuordnung. Der Meilensteinplan zeigt die wichtigsten Zwischenziele und deren geplante Fertigstellung.}

% --------------------------------------------------------
% Terminplanung und Meilensteine
% --------------------------------------------------------
\subsection*{Meilensteinplan}

\placeholderbox{Beispiel}{%
\begin{table}[H]
\centering
\begin{tabularx}{\linewidth}{lXl}
\hline
\textbf{Meilenstein} & \textbf{Beschreibung} & \textbf{Termin} \\
\hline
Requirement-Analyse & Anforderungen definieren und dokumentieren & KW 5 \\
Konzeptphase & Verschiedene Lösungsansätze entwickeln und bewerten & KW 7 \\
Implementierung Phase 1 & Kernfunktionalitäten umsetzen & KW 12 \\
Implementierung Phase 2 & Optimierung und erweiterte Features & KW 14 \\
Prototyp \& Testing & Tests durchführen und Ergebnisse dokumentieren & KW 16 \\
Dokumentation & Diplomarbeit finalisieren & KW 17 \\
Abgabe & Finale Abgabe & KW 18 \\
\hline
\end{tabularx}
\caption{Meilensteinplan der Diplomarbeit}
\label{tab:meilensteinplan}
\end{table}
}

% --------------------------------------------------------
% Bewertung der Zielerreichung
% --------------------------------------------------------
\section{Evaluierung}

\emph{Die Evaluierung dokumentiert, wie die gesetzten Ziele erreicht wurden und welche Anpassungen während der Projektlaufzeit vorgenommen wurden. Sie bietet eine Reflexion über die Projektabwicklung und zeigt eventuelle Abweichungen vom ursprünglichen Plan.}

% --------------------------------------------------------
% Arbeitszeiten und Stundenerfassung Autor A
% --------------------------------------------------------
\section{Zeiterfassung \htlAuthorNameA}

\emph{Die Zeiterfassung dokumentiert die investierte Arbeitszeit pro Projektphase und Woche. Dies ermöglicht eine Nachverfolgung der Ressourcennutzung und dient als Grundlage für zukünftige Zeitschätzungen.}

\placeholderbox{Beispiel}{%
\begin{table}[H]
\centering
\renewcommand{\arraystretch}{1.2}
\begin{tabularx}{\textwidth}{l r r X}
\textbf{KW} & \multicolumn{2}{c}{\textbf{Arbeitszeit [h]}} & \textbf{Aufgabenbeschreibung} \\
 & Schule & Freizeit & \\
\hline
KW22 & 10 & 2 &Anforderungsanalyse, Meetings mit Betreuung \\
KW23 & 10 & 4 & Recherche State of the Art, Literature Review \\
\ldots & \ldots & \ldots & \ldots \\
\hline
\textbf{Summe} & \textbf{204} & \textbf{185} &  \\
\hline
\hline
\end{tabularx}
\caption{Wöchentliche Zeiterfassung \htlAuthorNameB}
\label{tab:zeiterfassung_a}
\end{table}
}

% --------------------------------------------------------
% Unterscrift A
% --------------------------------------------------------
\begin{tabularx}{\textwidth}{l p{1cm} l p{1cm} X}

\htlOrt, \todayshort & & \htlAuthorNameA & & \hrulefill \\
\emph{Ort, Datum} & & & & \emph{Unterschrift} \vspace{2cm}\\ 

\end{tabularx}


% --------------------------------------------------------
% Arbeitszeiten und Stundenerfassung Autor B
% --------------------------------------------------------
\section{Zeiterfassung \htlAuthorNameB}


\emph{Die Zeiterfassung dokumentiert die investierte Arbeitszeit pro Projektphase und Woche. Dies ermöglicht eine Nachverfolgung der Ressourcennutzung und dient als Grundlage für zukünftige Zeitschätzungen.}

\placeholderbox{Beispiel}{%
\begin{table}[H]
\centering
\renewcommand{\arraystretch}{1.2}
\begin{tabularx}{\textwidth}{l r r X}
\textbf{KW} & \multicolumn{2}{c}{\textbf{Arbeitszeit [h]}} & \textbf{Aufgabenbeschreibung} \\
 & Schule & Freizeit & \\
\hline
KW22 & 10 & 2 &Anforderungsanalyse, Meetings mit Betreuung \\
KW23 & 10 & 4 & Recherche State of the Art, Literature Review \\
\ldots & \ldots & \ldots & \ldots \\
\hline
\textbf{Summe} & \textbf{204} & \textbf{185} &  \\
\hline
\hline
\end{tabularx}
\caption{Wöchentliche Zeiterfassung \htlAuthorNameB}
\label{tab:zeiterfassung_b}
\end{table}
}


% --------------------------------------------------------
% Unterscrift B
% --------------------------------------------------------

\begin{tabularx}{\textwidth}{l p{1cm} l p{1cm} X}

\htlOrt, \todayshort & & \htlAuthorNameB & & \hrulefill \\
\emph{Ort, Datum} & & & & \emph{Unterschrift} \vspace{2cm}\\ 

\end{tabularx}